\section{Introduction}
\label{sec:introduction}

The success of deep learning in computer vision relies on the assumption that training and testing data share the same distribution. In practice, this assumption is frequently violated: models trained in one environment often fail when deployed in another due to domain shifts in factors such as lighting, viewpoint, or sensor characteristics. Domain Generalization (DG) addresses this challenge by learning models from multiple source domains that generalize robustly to unseen target domains~\citep{zhou2022domain}. Existing DG paradigms include data augmentation to synthesize diversity, meta-learning to acquire generalization strategies, and domain alignment, which explicitly learns domain-invariant feature representations.

Within the domain alignment paradigm, the goal is to minimize discrepancy between feature distributions of different source domains. Prior methods have employed statistical metrics such as Maximum Mean Discrepancy (MMD)~\citep{li2018domain_mmd} or adversarial training inspired by Domain-Adversarial Neural Networks (DANN)~\citep{ganin2016domain}. However, adversarial approaches often suffer from training instability and mode collapse, while MMD-based methods are sensitive to kernel selection and bandwidth hyperparameters. We argue that Optimal Transport (OT) offers a more principled alternative. OT provides a geometrically grounded framework for comparing probability distributions, yielding meaningful distances even when distribution supports do not overlap, which is a common scenario in DG. The Sinkhorn algorithm~\citep{cuturi2013sinkhorn} enables efficient, differentiable OT computation suitable for deep learning.

Despite its theoretical appeal, a critical limitation has hindered OT's effectiveness in multi-source DG: the instability of the alignment pivot. Aligning multiple source domains requires a common reference distribution. Conventional approaches use batch-derived targets, such as features from a randomly selected source domain or the pooled features within the current minibatch. These ephemeral pivots exhibit high variance, fluctuating substantially with each training step due to stochastic data sampling. This \textit{pivot instability} creates an unreliable optimization landscape, introducing noisy and often conflicting gradient signals that hinder convergence to domain-invariant representations.

\begin{sloppypar} 
This paper addresses this foundational challenge through SODA (\underline{S}tabilized \underline{O}ptimal Transport Based \underline{D}omain \underline{A}lignment), a framework that stabilizes OT-based alignment via a consistent target distribution. Our key contribution is the Dynamic Feature Reservoir (DFR), a high-capacity, first-in-first-out (FIFO) feature queue that serves as a non-parametric pivot distribution. The large capacity buffers the pivot from minibatch fluctuations, ensuring stability, while the FIFO update rule enables gradual adaptation as the feature extractor improves during training. This design yields a pivot that is stable over short timescales yet adaptive over long timescales. By aligning all source domains to this reservoir using a regularized OT loss, SODA provides a smooth optimization landscape for learning robust features. Extensive experiments on seven DG benchmarks, including five real-world datasets, demonstrate that SODA achieves consistently competitive or state-of-the-art accuracy with a lower variance of the alignment gradient than batch-derived pivots.
\end{sloppypar} 

The DFR design draws inspiration from memory queues in self-supervised learning (MoCo~\citep{he2020momentum}) and semi-supervised learning (Mean Teacher~\citep{tarvainen2017mean}). However, the DFR serves a fundamentally different mathematical role: rather than providing negative samples for contrastive learning or enforcing temporal consistency through parameter averaging, the DFR acts as a Monte Carlo estimator for the target distribution in Wasserstein space. This enables computing transport plans against a sample size $K$ that is substantially larger than the mini-batch size $B$, addressing the sample complexity limitations inherent in high-dimensional distribution matching~\citep{dong2020how}. A detailed comparison of these mechanisms is provided in Section~\ref{sec:stability_mechanisms}.

Our main contributions are:
\begin{itemize}
    \item We identify that standard mini-batch OT alignment suffers from high gradient variance due to \textit{pivot instability}, where batch-derived target distributions are statistically impoverished and fluctuate across iterations.
    \item We propose the Dynamic Feature Reservoir (DFR), which decouples the alignment target's sample size ($K$) from the training batch size ($B$), enabling low-variance geometric alignment without increasing batch memory requirements.
    \item We demonstrate that stabilizing the transport target yields robust generalization, achieving consistently competitive or state-of-the-art accuracy across seven DG benchmarks while reducing the pivot-induced variance of the alignment gradient compared to batch-derived alignment methods.
\end{itemize}